\subsection[Selection Sort]{Selection Sort \cite{alnihoud2010enhancement}}
    \subsubsection{Overview}
        \begin{itemize}
            \item \textbf{History and Origin:} A simple and intuitive in-place comparison-based sorting algorithm.
            \item \textbf{Efficiency:} Generally performs worse than similar algorithms like Insertion Sort. Best suited for situations where memory writes are costly.
            \item \textbf{Application:} Useful for small datasets or when auxiliary memory is extremely limited.
        \end{itemize}
    \subsubsection{Core Concepts}
        \begin{itemize}
            \item The algorithm divides the input list into two parts: a sorted sublist built up from left to right at the front, and a sublist of the remaining unsorted items.
            \item Initially, the sorted sublist is empty and the unsorted sublist is the entire list.
            \item It proceeds by finding the smallest element in the unsorted sublist, swapping it with the leftmost unsorted element, and moving the sublist boundary one element to the right.
        \end{itemize}

    \subsubsection{Step-by-step Explanation}
        \begin{itemize}
            \item Consider an array of n elements, a[0\ldots (n - 1)].
            \item \textbf{Step 1:} Start a loop with variable \textit{i} from 0 to n - 2.
            \item \textbf{Step 2:} Assume the element at \textit{i} is the minimum (\textit{min} = \textit{i}).
            \item \textbf{Step 3:} Iterate through the remaining array from \textit{j} = \textit{i} + 1 to n - 1 to find the actual minimum element.
            \begin{itemize}
                \item If a[j] < a[min], update \textit{min} = \textit{j}.
            \end{itemize}
            \item \textbf{Step 4:} Swap a[i] with a[min].
            \item \textbf{Step 5:} Increase \textit{i} by 1 and repeat until the array is fully sorted.
        \end{itemize}

    \subsubsection{Pseudo Code}
        \begin{itemize}
            \item \textbf{Algorithm:} Selection Sort
            \item \textbf{Input:} An array A of n elements
        \end{itemize}
        \begin{verbatim}
        1. i = 0
        2. While i < n - 1:
        3.    min = i
        4.    For j = i + 1 to n - 1:
        5.       If A[min] > A[j]:
        6.          min = j
        7.    Swap(A[i], A[min])
        8.    i = i + 1
        \end{verbatim}
    \subsubsection{Complexity Analysis}
        \begin{itemize}
            \item \textbf{Time Complexity:}
            \begin{itemize}
                \item Best case: $O(n^2)$
                \item Average case: $O(n^2)$
                \item Worst case: $O(n^2)$ 
                    \subitem The algorithm scans the entire remaining array regardless of the initial order.
            \end{itemize}
            \item \textbf{Space Complexity:} $O(1)$ 
                \subitem In-place algorithm, requiring only a constant amount of extra memory for swapping.
        \end{itemize}
   \subsubsection{Variants and Optimizations}
        \begin{itemize}
            \item \textbf{Current Implementation (Standard):} 
            \begin{itemize}
                \item The provided C++ code implements the standard unidirectional Selection Sort, which scans for the minimum element from left to right in a single pass.
            \end{itemize}
            \item \textbf{Theoretical Variant: Bidirectional Selection Sort \cite{alnihoud2010enhancement}:}
            \begin{itemize}
                \item \textbf{Problem:} Conventional Selection Sort only finds the minimum element in each pass, leaving the rest of the array unsorted.
                \item \textbf{Solution:} In each pass, simultaneously find both the minimum and maximum elements, placing them at the beginning and end of the unsorted section respectively.
                \item \textbf{Benefit:} Reduces the total number of iterations by half compared to the implemented version, improving practical execution time.
            \end{itemize}
        \end{itemize}